% Generated draft: requires mathematical and source review.
\documentclass[11pt]{article}
\usepackage[margin=25mm]{geometry}
\usepackage{fontspec}
\setmainfont{Latin Modern Roman}
\usepackage{amsmath,amssymb,mathtools,unicode-math}
\setmathfont{Latin Modern Math}
\usepackage{enumitem,needspace,xurl,hyperref}
\setlength{\parindent}{0pt}
\setlength{\parskip}{5pt}
\setlength{\emergencystretch}{3em}
\newcommand{\N}{\mathbb{N}}
\newcommand{\Z}{\mathbb{Z}}
\newcommand{\Q}{\mathbb{Q}}
\newcommand{\R}{\mathbb{R}}
\newcommand{\C}{\mathbb{C}}
\newcommand{\F}{\mathbb{F}}
\newcommand{\A}{\mathbb{A}}
\newcommand{\PP}{\mathbb{P}}
\newcommand{\E}{\mathbb{E}}
\newcommand{\eps}{\varepsilon}
\newcommand{\dd}{\,\mathrm{d}}
\begin{document}
\section*{UnsolvedMath ID: 801089}
\textbf{Source problem code:} catalogue entry
\subsection{Definitions and mathematical statement}
Let \(\operatorname{Aut}_{\mathrm{poly}}(\C^3)\) denote the group of polynomial maps \(\C^3\to\C^3\) having polynomial inverses. Define \(\mathrm{Tame}(3)\) to be its subgroup generated by affine automorphisms \(z\mapsto Az+b\), where \(A\in\operatorname{GL}_3(\C)\) and \(b\in\C^3\), and triangular automorphisms. Using an upper triangular coordinate convention, the latter have the form
\[
(z_1,z_2,z_3)\longmapsto
\bigl(a_1z_1+P_1(z_2,z_3),\ a_2z_2+P_2(z_3),\ a_3z_3+c\bigr),
\]
where \(a_1,a_2,a_3\in\C\setminus\{0\}\), \(c\in\C\), \(P_1\in\C[z_2,z_3]\), and \(P_2\in\C[z_3]\).

Let \(\mathcal C\) be the particular \(\mathrm{CAT}(0)\) complex attributed in the supplied source to Lamy and Przytycky [LP], equipped with their isometric action of \(\mathrm{Tame}(3)\). Here \(\mathrm{CAT}(0)\) means that the geodesic metric satisfies the Euclidean comparison inequality for geodesic triangles. The extracted record does not specify the construction, metric, or action on \(\mathcal C\); the question refers to that particular construction, rather than to an arbitrary isometric action on a \(\mathrm{CAT}(0)\) space.

Under the conventional interpretation of \(\lambda(f)\) as the first dynamical degree, put
\[
\deg(g)=\max_{1\leq i\leq3}\deg(g_i),
\qquad
\lambda(f)=\lim_{n\to\infty}\bigl(\deg(f^n)\bigr)^{1/n},
\]
where coordinate degrees are total polynomial degrees and \(f^n\) is the \(n\)-fold iterate. The extracted record does not define \(\lambda\), so this interpretation requires confirmation against the complete cited source.

\textbf{Question.} Characterize precisely the elements
\[
\bigl\{f\in\mathrm{Tame}(3):\lambda(f)=1\bigr\}
\]
in terms of the isometries they induce on \(\mathcal C\). Thus the requested characterization should give necessary and sufficient conditions expressed through this action. The supplied question does not propose a specific condition or a particular classification of isometries to use.

For context, the source attributes to Shestakov and Umirbaev [SU] the theorem that \(\mathrm{Tame}(3)\) is a proper subgroup of \(\operatorname{Aut}_{\mathrm{poly}}(\C^3)\), contrasting this with equality of the analogous groups in dimension two. These are supplied background assertions; this draft does not independently verify them.
\subsection{Short English statement}
For a tame polynomial automorphism of \(\C^3\), characterize when its dynamical degree equals one using its action on the particular \(\mathrm{CAT}(0)\) complex attributed to Lamy and Przytycky. The meaning of \(\lambda\) and the precise complex and action require confirmation from the cited source. The registry labels the question open; this is a statement-only draft requiring human mathematical review.
\subsection{Sources}
\begin{itemize}
\item Local UnsolvedMath registry, permanent ID \(801089\), exact source problem code \texttt{catalogue entry}, title ``Henon Maps: Open Problems --- Question 41.'' The record attributes the problem to Berger and Bedford, proposed in 2023, and supplies the status \texttt{open}; this status is not an independently verified literature assessment.
\item Berger and Bedford, \emph{Henon Maps: A List of Open Problems} (2023), Question 41. Supplied source URL: \url{https://arxiv.org/abs/2312.03907}. The record reports access on 2026-10-16 and extraction by \texttt{arxiv\_latex}; these provenance statements are supplied metadata, not a claim of access or verification by this draft.
\item Reference label [SU]: Shestakov and Umirbaev, cited in the extracted question for the properness of \(\mathrm{Tame}(3)\). No title, date, URL, or full bibliographic entry is supplied for [SU].
\item Reference label [LP]: Lamy and Przytycky, as named in the extracted question, cited for the \(\mathrm{CAT}(0)\) complex and isometric action. No title, date, URL, full bibliographic entry, or construction is supplied for [LP].
\end{itemize}
\end{document}
